% Part: incompleteness
% Chapter: representability-in-q
% Section: c-representable

\documentclass[../../../include/open-logic-section]{subfiles}

\begin{document}

\olfileid[hi]{inc}{req}{cre}
\olsection{$C$ के फलन $\Th{Q}$ में निरूपणीय हैं}

हमें दिखाना है कि $C$ का प्रत्येक फलन $\Th{Q}$ में निरूपणीय है। अंततः
हमें दिखाना होगा कि $C$ के प्रत्येक $k$-स्थानीय फलन
$f(x_0,\dots,x_{k-1})$ को उसका निरूपण करने वाला कोई सूत्र
$!A_f(x_0,\dots,x_{k-1},y)$ कैसे नियत किया जाए।

\begin{lem}
प्राकृतिक संख्याएँ $n$ और $m$ दी गई हों। यदि $n \neq m$, तो
$Q \Proves \num n \neq \num m$।
\end{lem}

\begin{proof}
$n$ पर आगमन का उपयोग करके दिखाएँ कि प्रत्येक $m$ के लिए, यदि $n \neq m$,
तो $Q \Proves \eq/[\num n][\num m]$।

आधार मामले में $n = 0$। यदि $m$,~$0$ के बराबर नहीं है, तो किसी प्राकृतिक
संख्या $k$ के लिए $m = k + 1$। हमारे पास अभिगृहीत
$\lforall[x][\eq/[0][x']]$ है। किसी परिमाणक-अभिगृहीत से $x$ के स्थान पर
$\num k$ रखकर हम $\eq/[0][\num k']$ निष्कर्षित कर सकते हैं। किंतु
$\num k'$ केवल $\num m$ है।

आगमन चरण में हम मान सकते हैं कि दावा $n$ के लिए सत्य है और $n+1$ पर विचार
करते हैं। $m$ कोई भी प्राकृतिक संख्या हो। दो संभावनाएँ हैं: या तो $m = 0$,
अथवा किसी $k$ के लिए $m = k+1$। पहला मामला ऊपर की तरह हल होता है। दूसरे
मामले में मान लें $n+1 \neq k+1$। तब $n \neq k$। $n$ के लिए आगमन
परिकल्पना से $\Th{Q} \Proves \eq/[\num n][\num k]$। हमारे पास अभिगृहीत
$\lforall[x][\lforall[y][\eq[x'][y'] \lif \eq[x][y]]]$ है। किसी
परिमाणक-अभिगृहीत से हमारे पास
$\eq[\num n'][\num k'] \lif \eq[\num n][\num k]$ है। प्रतिज्ञप्तिक तर्क
का उपयोग करके हम $\Th{Q}$ में
$\eq/[\num n][\num k] \lif \eq/[\num n'][\num k']$ निष्कर्षित कर सकते
हैं। मोडस पोनेन्स से हमें $\eq/[\num n'][\num k']$ मिलता है, और यही हमें
चाहिए, क्योंकि $\num k'$,~$\num m$ है।
\end{proof}

ध्यान दें कि उपपत्ति बहुत बड़ा दावा नहीं करती: मूलतः वह कहती है कि
$\Th{Q}$ सिद्ध कर सकता है कि अलग-अलग अंक-सूचक अलग वस्तुओं को निर्दिष्ट
करते हैं। उदाहरणतः $\Th{Q}$ सिद्ध करता है कि $0'' \neq 0'''$। किंतु इसे
सामान्य रूप में दिखाने के लिए कुछ सावधानी चाहिए। यह भी ध्यान दें कि यद्यपि
हम आगमन का उपयोग कर रहे हैं, यह $\Th{Q}$ के \emph{बाहर} का आगमन है।

हम शून्य, उत्तराधिकारी, जोड़, गुणन, समता के अभिलाक्षणिक फलन और प्रक्षेपों
का निरूपण कर सकेंगे। प्रत्येक मामले में उपयुक्त निरूपक सूत्र पूरी तरह सीधा
है; उदाहरणतः शून्य का निरूपण सूत्र
\[
y = 0,
\]
उत्तराधिकारी का निरूपण सूत्र
\[
x_0' = y,
\]
और जोड़ का निरूपण सूत्र
\[
x_0 + x_1 = y.
\]
द्वारा होता है। काम यह दिखाने में है कि $\Th{Q}$ संबद्ध कथनों को सिद्ध कर
सकता है; उदाहरणतः ऊपर के सूत्र द्वारा जोड़ का निरूपण होने का अर्थ यह
दिखाना है कि प्राकृतिक संख्याओं $m$ और $n$ के प्रत्येक युग्म के लिए
$\Th{Q}$ सिद्ध करता है
\[
\eq[\num n + \num m][\num {n+m}]
\]
और
\[
\lforall[y][(\eq[(\num n + \num m)][y] \lif \eq[y][\num {n+m}])].
\]

संयोजन के बारे में क्या? मान लें $h$ इस प्रकार परिभाषित है:
\[
h(x_0,\dots,x_{l-1}) = f(g_0(x_0,\dots,x_{l-1}), \dots,
g_{k-1}(x_0,\dots,x_{l-1})).
\]
और हम पहले ही ऐसे सूत्र $!A_f, !A_{g_0}, \dots, !A_{g_{k-1}}$ पा चुके हैं
जो क्रमशः फलनों $f,g_0,\dots,g_{k-1}$ का निरूपण करते हैं। तब हम~$h$ का
निरूपण करने वाला सूत्र $!A_h$ इस प्रकार परिभाषित कर सकते हैं:
$!A_h(x_0,\dots,x_{l-1},y)$ हो
\begin{multline*}
  \lexists[z_0\dots][\lexists[z_{k-1}][(!A_{g_0}(x_0,\dots,x_{l-1},z_0) \land
  \dots \land !A_{g_{k-1}}(x_0,\dots,x_{l-1},z_{k-1}) \land]]\\
  {}!A_f(z_0,\dots,z_{k-1},y)).
\end{multline*}

अंततः अपरिबद्ध खोज पर विचार करें। मान लें $g(x,\vec z)$ नियमित है और
$\Th{Q}$ में, मान लें सूत्र $!A_g(x,\vec z,y)$ द्वारा, निरूपणीय है। $f$
को $f(\vec z) = \mu x \; g(x,\vec z)$ द्वारा परिभाषित करें। हम~$f$ का
निरूपण करने वाला कोई सूत्र $!A_f(\vec z,y)$ खोजना चाहते हैं। एक स्वाभाविक
चुनाव है:
\[
!A_f(\vec z,y) \ident !A_g(y,\vec z,0) \land \lforall[w][(w < y
\lif \lnot !A_g(w,\vec z,0))].
\]

कुछ अतिरिक्त उपपत्तियों का उपयोग करके दिखाया जा सकता है कि यह काम करता
है; उदाहरणतः:

\begin{lem}
  प्रत्येक चर $x$ और प्रत्येक प्राकृतिक संख्या $n$ के लिए $\Th{Q}$,
  $\eq[(x' + \num n)][(x + \num n)']$ को सिद्ध करता है।
\end{lem}

यह फिर उल्लेखनीय है कि यह उस कथन से दुर्बल है कि $\Th{Q}$,
$\lforall[x][\lforall[y][(x' + y) = (x +y)']]$ को सिद्ध करता है (जो असत्य
है)।

\begin{proof}
सामान्यतः प्रमाण $n$ पर आगमन द्वारा है। आधार मामले में $n = 0$; हमें
दिखाना है कि $\Th{Q}$,~$\eq[(x'+0)][(x+0)']$ को सिद्ध करता है। किंतु हमारे
पास है:
\begin{eqnarray*}
& & \eq[(x' + 0)][x'] \quad \text{अभिगृहीत 4 से} \\
& & \eq[(x + 0)][x] \quad \text{अभिगृहीत 4 से} \\
& & \eq[(x+0)'][x'] \quad \text{पंक्ति 2 से} \\
& & \eq[(x' + 0)][(x + 0)'] \quad \text{पंक्तियों 1 और 3 से।}
\end{eqnarray*}
आगमन चरण में हम मान सकते हैं कि हमने $\Th{Q}$ में
$\eq[(x' + \num n)][(x + \num n)']$ व्युत्पन्न किया है। चूँकि
$\num{n+1}$,~$\num n'$ है, हमें दिखाना है कि $\Th{Q}$,
$\eq[(x' + \num n')][(x + \num n')']$ को सिद्ध करता है। समताओं की निम्न
शृंखला $\Th{Q}$ में व्युत्पन्न की जा सकती है:
\begin{eqnarray*}
(x' + \num n') & \eq & (x' + \num n)' \quad \text{अभिगृहीत 5} \\
& \eq & (x + \num n')' \quad \text{आगमन परिकल्पना से।}
\end{eqnarray*}
\end{proof}

\begin{lem}
\begin{enumerate}
\item $\Th{Q}$,~$\lnot (x < \num 0)$ को सिद्ध करता है।
\item प्रत्येक प्राकृतिक संख्या $n$ के लिए $\Th{Q}$ सिद्ध करता है
\[
x < \num {n+1} \lif (\eq[x][\num 0] \lor \dots \lor \eq[x][\num n]).
\]
\end{enumerate}
\end{lem}

\begin{proof}
आइए 1 और 2 का कुछ भाग अनौपचारिक रूप से करें (अर्थात् औपचारिक
!!{derivation} बनाने के केवल संकेत दें)।

भाग 1 के लिए $<$ की परिभाषा से हमें $\Th{Q}$ में
$\lnot \lexists[y][\eq[(y'+x)][0]]$ सिद्ध करना है, जो प्रथम-कोटि
तर्कशास्त्र के अभिगृहीतों और नियमों से
$\lforall[y][\eq/[(y' + x)][0]]$ के समतुल्य है। विचार यह है: मान लें
$\eq[(y'+x)][0]$। यदि $x$,~0 है, तो हमारे पास
$\eq[(y' + 0)][0]$ है। किंतु $\Th{Q}$ के अभिगृहीत 4 से
$\eq[(y' + 0)][y']$, और अभिगृहीत 2 से $\eq/[y'][0]$, जो विरोधाभास है।
अतः $\lforall[y][\eq/[(y' +x)][0]]$। यदि $x$,~$0$ नहीं है, तो अभिगृहीत 3
से कोई $z$ ऐसा है कि $\eq[x][z']$। किंतु तब हमारे पास
$\eq[(y' + z')][0]$ है। अभिगृहीत~5 से $\eq[(y' + z)'][0]$, जो फिर
अभिगृहीत~2 के विरुद्ध है।

भाग 2 के लिए $n$ पर आगमन का उपयोग करें। आधार मामला $n =0$ लें। उस मामले
में हमें $x < \num 1 \lif \eq[x][\num 0]$ दिखाना है। मान लें
$x < \num 1$। तब $<$ के परिभाषक अभिगृहीत से हमारे पास
$\lexists[y][\eq[(y'+x)][0']]$ है। मान लें $y$ में यह गुण है; अतः हमारे
पास $\eq[y'+x][0']$ है।

हमें $\eq[x][0]$ दिखाना है। अभिगृहीत 3 से यदि $x$,~0 नहीं है, तो वह
किसी $z$ के लिए $z'$ के बराबर है। तब हमारे पास
$\eq[(y' + z')][0']$ है। $\Th{Q}$ के अभिगृहीत~5 से
$\eq[(y'+z)'][0']$। अभिगृहीत 1 से $\eq[(y' + z)][0]$। किंतु परिभाषा से
इसका अर्थ $z < 0$ है, जो भाग~1 के विरुद्ध है।
\end{proof}

\begin{explain}
हमने दिखाया है कि संगणनीय फलनों के समुच्चय को $\Th{Q}$ में निरूपणीय फलनों
के समुच्चय के रूप में अभिलक्षित किया जा सकता है। वास्तव में प्रमाण अधिक
व्यापक है। निरूपणीयता की परिभाषा से यह देखना कठिन नहीं है कि $\Th{Q}$ का
कोई भी विस्तार (अथवा कोई ऐसा सिद्धांत जिसमें $\Th{Q}$ की व्याख्या की जा
सके) संगणनीय फलनों का निरूपण कर सकता है; किंतु विलोमतः, जिस किसी
!!{derivation}-तंत्र में प्रमाण की धारणा संगणनीय हो, उसमें प्रत्येक निरूपणीय
फलन संगणनीय है। अतः उदाहरण के लिए संगणनीय फलनों के समुच्चय को पियानो
अंकगणित, बल्कि ज़र्मेलो--फ्रेंकेल समुच्चय-सिद्धांत में भी निरूपित फलनों के
समुच्चय के रूप में अभिलक्षित किया जा सकता है। जैसा ग्योडेल (G\"odel) ने
उल्लेख किया, यह कुछ आश्चर्यजनक है। हम देखेंगे कि सिद्धता के प्रश्न इस बात
के प्रति अत्यंत संवेदनशील होते हैं कि हम किस सिद्धांत पर विचार कर रहे हैं;
मोटे तौर पर अभिगृहीत जितने प्रबल होंगे, उतना अधिक सिद्ध किया जा सकेगा।
किंतु अभिगृहीती सिद्धांतों के व्यापक वर्ग में निरूपणीय फलन ठीक संगणनीय फलन
ही हैं।
\end{explain}

\end{document}
